%%%%%
%@TheDoctorRAB
%%%%%
%
%%%%%
%standard report format for journal papers
%%%%%
%
%%%%%
%PDF tagging accessible file
%compile with pdflatex
%run $show-pdf-tags --xml filename.pdf
%%%%%
%
%%%%%
%REFERENCES
%neup.bst - numbered citations in order of appearance, short author list with et al in reference section
%nsf.bst - numbered citations in order of appearance, full author list in references section
%standard.bst - citations with author last name with et al for more than 2 authors; full author list in references section
%ans.bst is for ANS only. 
%
%author = {Lastname, Firstname and Lastname, Firstname and Lastname, Firstname} for all bst formats
%bst renders the author list itself
%
%author = {{Nuclear Regulatory Commission}} if the author is an organization, institution, etc., and not people
%
%title = {{}} for all
%
%for all - use \citep{-} - [1] or (Borrelli, 2021) in the text
%standard.bst \cite{-} - Borrelli (2021) in the text
%standard.bst lists references alphabetically
%the rest list numerically
%%%%%


%%%%% enable tagging for accessibility
\DocumentMetadata{
    tagging=on,
    check-tagging-status,
    pdfversion=2.0,
    pdfstandard=ua-2,
    lang=en-US,
    debug={log=all},
    testphase={latest,tagpdf},
    tagging-setup={math/setup={mathml-AF,mathml-SE}}
}
\tagpdfsetup{math/alt/use}
%%%%%


%%%%% general 
\documentclass[11pt,a4paper]{article}
\usepackage[lmargin=1in,rmargin=1in,tmargin=1in,bmargin=1in]{geometry}
\usepackage[pagewise]{lineno} %linenumbering restarts each page
%\usepackage[running]{lineno} %linenumbering continuous
\usepackage{mmap} %helps make pdfs copy and paste
\usepackage{setspace}
\usepackage{ulem} %strikethrough - do not \sout{\cite{}}
\usepackage[dvipsnames,table]{xcolor} %change font color (table loads colortbl)
\usepackage{graphicx}
\usepackage{tablefootnote}
\usepackage{footnotehyper}
\usepackage{float}
%\usepackage{mypythonhighlight,verbatim}
%\usepackage{subfig}
\usepackage[yyyymmdd]{datetime} %date format
\renewcommand{\dateseparator}{.}
%\graphicspath{{$TEXIMG/}} %path to graphics
\setcounter{secnumdepth}{5} %set subsection to nth level
\usepackage{needspace}
\usepackage[stable,hang,flushmargin]{footmisc} %footnotes in section titles and no indent; standard.bst
\usepackage[inline]{enumitem} %https://www.tug.org/texlive//Contents/live/texmf-dist/doc/latex-dev/latex-lab/latex-lab-enumitem.pdf
\usepackage{boldline}
\usepackage{makecell}
\usepackage{booktabs}
\usepackage{amssymb}
\usepackage{gensymb}
\usepackage{amsmath,nicefrac}
\usepackage{physics}
\usepackage{lscape}
\usepackage{array}
\usepackage{chngcntr}
\usepackage{sectsty}
\usepackage{textcomp}
\usepackage{lastpage}
\usepackage{xargs} %for \newcommandx
\usepackage[colorinlistoftodos,prependcaption,textsize=tiny]{todonotes} %makes colored boxes for commenting
\usepackage{soul}
\usepackage{color}
\usepackage{marginnote}
\usepackage[figure,table]{totalcount}
\usepackage{parskip}
%%%%%


%%%%% links and tags
\usepackage{hyperref}
\hypersetup{
    colorlinks,
    linkcolor=black,
    citecolor=black,
    urlcolor=blue,
    pdftitle={Homework},
    pdfauthor={R. A. Borrelli},
    pdfkeywords={homework},
    pdflanguage={en-US},
    pdfsubject={homework},
    pdfdisplaydoctitle=true 
} 
\usepackage[capitalise]{cleveref} %must be loaded after hyperref
%%%%%


%%%%% tikz
\usepackage{pgf}
\usepackage{tikz} % required for drawing custom shapes
\usetikzlibrary{shapes,arrows,automata,trees}
%%%%%


%%%%% fonts
\usepackage{times}
%arial - uncomment next two lines
%\usepackage{helvet}
%\renewcommand{\familydefault}{\sfdefault}
%%%%%


%%%%% references
%\usepackage[round,semicolon]{natbib} %for (Borrelli 2021; Clooney 2019) - standard.bst 
\usepackage[numbers,sort&compress]{natbib} %for [1-3] - nsf.bst, neup.bst
\setlength{\bibsep}{7pt} %sets space between references
\renewcommand{\bibsection}{\section*{\centering\bfseries\refname}} %bold, centered References heading
%\renewcommand{\bibsection}{} %suppresses large 'references' heading
%\renewcommand\bibpreamble{\vspace{\baselineskip}} %sets spacing after heading if not using default references heading
%%%%%


%%%%% tables and figures
\usepackage{longtable} %need to put label at top under caption then \\ - use spacing
\usepackage{tablefootnote}
\usepackage{tabularx}
\usepackage{multirow}
\usepackage{tabto} %general tabbed spacing
\usepackage{pdfpages}
\usepackage{wrapfig} %wraps figures around text
\setlength{\intextsep}{0.00mm}
\setlength{\columnsep}{1.00mm}
\usepackage[singlelinecheck=false,labelfont=bf]{caption}
\usepackage{subcaption}
\captionsetup[table]{justification=centering,skip=5pt,labelformat={default},labelsep=period,name={Table}} %sets a space after table caption
\captionsetup[figure]{justification=justified,skip=5pt,labelformat={default},labelsep=period,name={Figure}} %sets space above caption, 'figure' format
\captionsetup[wrapfigure]{justification=centering,aboveskip=0pt,belowskip=0pt,labelformat={default},labelsep=period,name={Fig.}} %sets space above caption, 'figure' format
\captionsetup[wraptable]{justification=centering,aboveskip=0pt,belowskip=0pt,labelformat={default},labelsep=period,name={Table}} %sets space above caption, 'figure' format
%%%%%


%%%%% watermark
%\usepackage[firstpage,vpos=0.63\paperheight]{draftwatermark}
%\SetWatermarkText{\shortstack{DRAFT\\do not distribute}}
%\SetWatermarkScale{0.20}
%%%%%


%%%%% cross referencing files
%\usepackage{xr} %for revisions - will cross reference from one file to here
%\externaldocument{/path/to/auxfilename} %aux file needed
%%%%%


%%%%% toc and glossaries
\usepackage[toc,title]{appendix}
\usepackage[acronym,nomain,nonumberlist]{glossaries}
\makenoidxglossaries
%do not load titlesec,titletoc
%%%%%


%%%%% editing
\newcommand{\edit}[1]{\textcolor{blue}{#1}} %shortcut for changing font color on revised text
\newcommand{\fn}[1]{\footnote{#1}} %shortcut for footnote tag
\newcommand*\sq{\mathbin{\vcenter{\hbox{\rule{.3ex}{.3ex}}}}} %makes a small square as a separator $\sq$
%\newcommand{\sk}[1]{\sout{#1}} %shortcut for default strikethrough - do not sk through citep
\newcommand\sk{\bgroup\markoverwith{\textcolor{red}{\rule[0.5ex]{1pt}{1pt}}}\ULon} %strikethrough with red line; not in \section{}
%\st{} does strikethrough using soul package but does not like acronyms
\newcommand{\blucell}{\cellcolor{aliceblue}} %use to shade in table cell
\newcommand{\grycekk}{\cellcolor{lightgray}} %use to shade in table cell
\newcommand{\whicell}{\cellcolor{antiquewhite}} %use to shade in table cell
%%%%%


%%%%% colors
%http://latexcolor.com/
%https://en.wikibooks.org/wiki/LaTeX/Colors#:~:text=black%2C%20blue%2C%20brown%2C%20cyan,be%20available%20on%20all%20systems.
\definecolor{aliceblue}{rgb}{0.94, 0.97, 1.0}
\definecolor{antiquewhite}{rgb}{0.98, 0.92, 0.84}
\definecolor{lightmauve}{rgb}{0.86, 0.82, 1.0}
\definecolor{brilliantlavender}{rgb}{0.96, 0.73, 1.0}
\definecolor{brandeisblue}{rgb}{0.0, 0.44, 1.0}
\definecolor{darkmidnightblue}{rgb}{0.0, 0.2, 0.4}

\newcommand{\x}{\cellcolor{aliceblue}} %use to shade in table cell
\newcommand{\y}{\cellcolor{lightgray}} %use to shade in table cell
\newcommand{\z}{\cellcolor{antiquewhite}} %use to shade in table cell
%%%%%


%%%%% acronyms
\newcommand{\acf}{\acrfull} %full acronym
\newcommand{\acl}{\acrlong} %long acronym
\newcommand{\acs}{\acrshort} %short acronym

\newcommand{\acfp}{\acrfullpl} %full acronym plural
\newcommand{\aclp}{\acrlongpl} %long acronym plural
\newcommand{\acsp}{\acrshortpl} %short acronym plural
%%%%%


%%%%% todonotes
\newcommandx{\cmt}[2][1=]{\todo[author=\textbf{STRUCTURE},tickmarkheight=0.15cm,linecolor=red,backgroundcolor=red!25,bordercolor=black,#1]{#2}}
\newcommandx{\con}[2][1=]{\todo[author=\textbf{CONTENT},tickmarkheight=0.15cm,linecolor=brilliantlavender,backgroundcolor=brilliantlavender,bordercolor=black,#1]{#2}}
\newcommandx{\rab}[2][1=]{\todo[noline,author=\textbf{RAB},backgroundcolor=Plum!25,bordercolor=black,#1]{#2}}


%\newcommandx{\jon}[2][1=]{\todo[noline,author=\textbf{ATTN: Johnson},backgroundcolor=blue!25,bordercolor=black,#1]{#2}}
%\newcommandx{\han}[2][1=]{\todo[noline,author=\textbf{ATTN: Haney},backgroundcolor=OliveGreen!25,bordercolor=black,#1]{#2}}
%\newcommandx{\rab}[2][1=]{\todo[author=\textbf{RAB},tickmarkheight=0.15cm,linecolor=Plum,backgroundcolor=Plum!25,bordercolor=black,#1]{#2}}
%\newcommandx{\han}[2][1=]{\todo[author=\textbf{ATTN: Haney},tickmarkheight=0.15cm,linecolor=OliveGreen,backgroundcolor=OliveGreen!25,bordercolor=OliveGreen,#1]{#2}}
%\newcommandx{\jon}[2][1=]{\todo[author=\textbf{ATTN: Johnson},tickmarkheight=0.15cm,linecolor=blue,backgroundcolor=blue!25,bordercolor=blue,#1]{#2}}


% highlighting 
\DeclareRobustCommand{\hlc}[1]{{\sethlcolor{LimeGreen}\hl{#1}}}
\makeatletter
    \if@todonotes@disabled
    \newcommand{\hlh}[2]{#1}
    \else
    \newcommand{\hlh}[2]{\han{#2}\hlc{#1}}
    \fi
    \makeatother

\DeclareRobustCommand{\hld}[1]{{\sethlcolor{CornflowerBlue}\hl{#1}}}
\makeatletter
    \if@todonotes@disabled
    \newcommand{\hlj}[2]{#1}
    \else
    \newcommand{\hlj}[2]{\jon{#2}\hld{#1}}
    \fi
    \makeatother

\DeclareRobustCommand{\hlf}[1]{{\sethlcolor{lightmauve}\hl{#1}}}
\makeatletter
    \if@todonotes@disabled
    \newcommand{\hlb}[2]{#1}
    \else
    \newcommand{\hlb}[2]{\rab{#2}\hlf{#1}}
    \fi
    \makeatother
%%%%%


%%%%% table alignments
\newcolumntype{L}[1]{>{\raggedright\let\newline\\\arraybackslash\hspace{0pt}}m{#1}} %uses \raggedright with m,p{} in table column
\newcolumntype{C}[1]{>{\centering\let\newline\\\arraybackslash\hspace{0pt}}m{#1}} %uses \raggedright with m,p{} in table column
\newcolumntype{R}[1]{>{\raggedleft\let\newline\\\arraybackslash\hspace{0pt}}m{#1}} %uses \raggedright with m,p{} in table column
%%%%%


%%%%% table contents
\makeatletter
\renewcommand\tableofcontents{%
    \@starttoc{toc}%
}
\makeatother

\makeatletter
\renewcommand\listoffigures{%
    \@starttoc{lof}%
}
\makeatother

\makeatletter
\renewcommand\listoftables{%
    \@starttoc{lot}%
}
\makeatother

\makeatletter
\newcommand*\ftp{\fontsize{16.5}{17.5}\selectfont}
\makeatother
%%%%%


%%%%% header and footer
\usepackage{fancyhdr}
\pagestyle{fancy}
\fancyhf{} %move page number to bottom right
%\renewcommand{\headrulewidth}{0pt} %turn off line in header
\lhead{\scriptsize NE 5350 — Nuclear Criticality Safety}
\chead{\scriptsize \textit{Homework 1}}
\rhead{\scriptsize \today}
\rfoot{\thepage}
%%%%%


%%%%% acronyms
\input{$ACRONYM/acronyms}
%\input{../acronyms/acronyms}
%%%%%


\begin{document}


\begin{titlepage}
    \title{
        NE 5350 -- Nuclear Criticality Safety\\
        Homework #1\\
    }
    \author{
        Gabriel Lewis
        \\ \\ \\
        University of Idaho $\sq$ Idaho Falls Center for Higher Education\\
        Nuclear Engineering and Industrial Management Department
        \\ \\ \\
        galewis99@outlook.com
    }
\clearpage %not have page number on title page
\maketitle
\vspace*{\fill}
\begin{flushright}{
        Total -- 100
}
\end{flushright}
\thispagestyle{empty} %start with page number 1 on second page
\end{titlepage}


\section{Mayak 1958 Criticality Accident}
\bold{Using the concepts of exponential power growth and reactivity feedback discussed in class,
describe the Mayak 1958 criticality accident.} \\

In a steady-state condition, in which a system remains in 'static reactivity', describes the state in which the 90$\%$ enriched uranyl nitrate solutoin remained in the cylinder, sub-critical as the reactivity was less than zero.

\begin{equation}
\rho = \frac{k_{eff}-1}{k_{eff}} < 0
\end{equation}

Reactivity ($\rho$) is dependent on the neutron population; as the neutron population increases within a system, $\rho$ increases, and the opposite is also true when $\rho$ decreases, the neutron population decreases. Within critical and sub-critical systems, the neutron population is proportional to power. \\

During operations after an experiment, the experimentalists were emptying the tank with six-liter bottles, being of favorable geometry. They began to deviate from procedure, as this work was tedious. After assuming the remaining quantity of uranyl nitrade were presumbably sub-critical, the experimentalists decided to unbolt and tilt the tank over to fill the bottles faster \cite{LA-13638}. As they tilt the tank over, as depicted in Figure 1, the solution inside favored one side of the tank and resembled the geometry of a sphere -reducing neutron leakage- and causing a prompt criticality event due to the exponential power growth of the surging sequential nuclear generations and reactivity feedback. \\

\begin{figure}[h!]
    \centering
    \includegraphics[width=0.8\linewidth]{HW1/Figures/cylindertoSpherlike.png}
    \label{fig-label-name}
    \caption{Solution Sitting Flat in Cylinder to Becoming Sphere-like}
\end{figure}

\subsection{Feedback Effects}
Once the favorable geometry (of the partially filled cylinderical tank) went out the window, which kept the solution distributed across a wide surface area, the uranyl nitrate solution went prompt super-critical. The reactivity of the system drastically increased due to the sudden change in geometry, exponentially increasing the neutron population in less than a second, and producing a bright flash of Cherenkov light. In this process, the solution explosively shot toward the ceiling, as much of the material changed phase to a gaseous form due to the sudden increase in power born from the criticality incident. This material change decreased reactivity, as the uranyl nitrate was no longer compactly contained in the tilted cylinder. The ejection of material fundamentally changed the system in a way that a re-criticality event was no longer possible. In the process of thermal expansion, prior to the ejection of fissile material, neutron leakage increased as reactivity decreased.   \\


\subsection{Exponential Growth}
Consider the production and "destruction" of neutrons and delayed neutron precursors (DNPs), which have been provided during the lecture, as an approximation of a zero-dimensional system:

\begin{equation}
\frac{d}{dt}n(t) = \frac{\rho(t) - \beta}{\Lambda} n(t) + \sum_{d=1}^{N_D} \lambda_{d}C_{d}(t)
\end{equation}
\begin{equation}
\frac{d}{dt}C_{d}(t) = \frac{\beta_{d}}{\Lambda} n(t) - \lambda_{d}C_{d}(t)
\end{equation}

where n(t) is the neutron population, $\rho$ is reactivity, $\beta_{d}$ is the effective delayed neutron fraction, $\lambda_{d}$ is the decay constant [s$^{-1}$], $\Lambda$ is the generation time, and $C_{d}$ is the DNP concentration. As reactivity ($\rho$) increases, due to a change of the geometry, so does the neutron population, that is proportional to power. In an analysis using OpenFOAM and Serpent2, Vega et al. determined that tilting the cylinder resulted in a reactivity change of $\Delta\rho_{geometry} = \$27.47$, with $k_{eff}$ ranging from 0.9142 to more than 1.17 prior to becoming subcritical after the material ejection \cite{Vega2015}. As discussed during Lecture 2, an insertion of $\$0.50$ produces exponential growth in the neutron population, let alone the insertion of $\$27.47$, and will continue to grow exponentially until the system fundamentally changes.

%\begin{equation}
%A(t) =
%\begin{pmatrix}
%\dfrac{\rho(t)-\beta}{\Lambda} & \lambda_1 & \lambda_2 & \cdots & \lambda_{N_D} \\[1.5ex]
%\dfrac{\beta_1}{\Lambda} & -\lambda_1 & & & \\[1.5ex]
%\dfrac{\beta_2}{\Lambda} & & -\lambda_2 & & \\[1.5ex]
%\vdots & & & \ddots & \\[1.5ex]
%\dfrac{\beta_{N_D}}{\Lambda} & & & & -\lambda_{N_D}
%\end{pmatrix}
%\end{equation}


%\begin{equation}
%\vec{x}(t) =
%\begin{pmatrix}
%n(t) \\
%C_1(t) \\
%\vdots \\
%C_{N_D}(t)
%\end{pmatrix}
%\end{equation}

%\begin{equation}
%\frac{d}{dt}\vec{x}(t) = A(t)\vec{x}(t)
%\end{equation}

%Which is only linear in the special case that A is constant. With the reactivity inserstion,

%\begin{equation}
%\vec{x}(t) = \vec{x_{o}} exp(A \cdot t)
%\end{equation}

%such that the exp is a matrix exponential. Therefore, the system will incur exponential growth until $\rho$ exceeds zero -becoming super critical - or the system fundamentally changes through reactivity feedback.

\newpage


\section{Demon Core Criticality Accidents}
\bold{In the first “Demon Core” criticality accident, the operator remained in the room after the
accident thereby significantly increasing the dose that he received. In the second “Demon
Core” criticality accident, the operator described physically separating the two hemi-spheres
though the power excursion had already been terminated by reactivity feedback effects. For
both of these criticality safety accidents, describe how the physics of reactivity feedback could have been used to modify the experiments to make them safer.}

With both criticality accidents, worker dose could have been prevented if the experiments were conducted remotely, like within a hot cell. However, these two accidents occured in the 1940s, so experimentalists were constrained by limitations in procedures to carry out criticality experiments or the necessary forsight to do so. Each accident occured after a slip in the fine motor skills of the worker. The following subsections will discuss how the experiments could have been made safer using reactivity feedback, despite the experimentalists being in the same room.

\subsection{Demon Core I}

The "demon core" was a 6.2 kg sphere of plutonium, composed of two hemispheres \cite{LA-13638}. The critical assembly was being created -- by hand -- by stacking 4.4 kg tungsten carbide bricks compactly around the sphere. The purpose of the tungsten carbide bricks was to reflect escaping neutrons back into the plutonium core, increasing reactivity. With the experimentalist, slipping his hand and dropping the final reflecting tungsten carbide brick on the assembly, he stayed in the room after the sphere went prompt supercritical. Before modifications to the experiment is explored, it should be stated that the experimentalist should have fled the scene to reduce the radiation dose exposure. \\
Tungsten carbide is dense, appropriate for a reflector material. As the bricks were arranged compactly around the core, it may have been safer for the distance to increase between the plutonium core and the surrounding tungsten carbide brcks. Not only would this give the experimentalist more time to flee in the case of an accident, but would increase the distance for the gamma and neutron radiation to travel, reducing the effective dose on the worker. \\
As the experiment composed of simple components, the sphere and the surrounding bricks, this setup could have been assembled in a warehouse with cranes to lift and lower the bricks and sphere into place. Once the system became critical, in this hypothetical excercise, the worker in this scenerio would still recieve a dose, but to a lesser extent.

\subsection{Demon Core II}

In the second demon core criticality accident, the core hemispheres were encased in beryllium as a reflector material. After the first man died the previous year, first and foremost, I would not be in the same room as the 6.2 kg demon core of plutonium. To modify this demonstration, the shims should have been relied on to perform the criticality experiment, as opposed to the screwdriver seperating the two hemispheres. The experimentalist did this to save on time, as he wanted to wrap things up before lunch. Futhermore, it should be stated that the demonstrator, after causing the prompt supercriticality, should have fled the room along with the others, rather than seperating the hemispheres and collecting a substantial, lethal dose. The hemispheres being left together may have brought upon re-criticalities. Recall that a decrease in temperature and thermal expansion increases reactvity, and vise-versa. \\

Assume that there is additional material surrounding the beryllium shell of the demon core. Within the hemisphers, suppose that there could be electric heaters in the material. Upon any transient conditions, the heaters would activate to increase the temperature within the hemisphers. Through the effect of Doppler Broadening, the resonance peaks would increase in the non-fissile nuclei of the shell, in which neutron absorption would be more likely to occur. This would fundamentally change the system, and drive down reactivity -- provided that the heaters would stay on. However, this may prove to be ineffective if the heater coils fail to raise temperature enough to induce this effect in the material.


\newpage







%%%%% header and footer
\pagestyle{fancyplain}
\fancyhf{} %clear header and footer 
\renewcommand{\headrulewidth}{0pt} %set line thickness in header; uncomment as is to remove line
\rfoot{\fancyplain{}{\thepage}}
\nolinenumbers
%%%%%


%%%%%%% citations
%
%%% make sure $BIBREF contains the path to each bib file in .bashrc
%%% make sure Overleaf has the latexmkrc pointing to where the bib and bst are located
%
\bibliographystyle{nsf}
\setlength{\bibhang}{0pt}
\bibliography{references}
%
%%%%%%%







%\begin{figure}[h!]
%    \centering
%    \includegraphics[width=0.X\textwidth]{filename}
%    \caption{Descriptive caption}
%    \label{fig-label-name}
%\end{figure}




%%%%% formatting
    \renewcommand{\thesection}{\Roman{section}}
%%%%%


%    \section{\acs{mcp} input decks} \label{app-mcnp}



\end{appendices}


\end{document}
